
\documentclass{article}
\usepackage{colortbl}
\usepackage{makecell}
\usepackage{multirow}
\usepackage{supertabular}

\begin{document}

\newcounter{utterance}

\centering \large Interaction Transcript for game `cladder', experiment `full\_v1.5\_default', episode 1743 with Dialogue Architects\_\_qwen3.5{-}9b\_\_SCoRe\_Qwen3.5{-}9B.
\vspace{24pt}

{ \footnotesize  \setcounter{utterance}{1}
\setlength{\tabcolsep}{0pt}
\begin{supertabular}{c@{$\;$}|p{.15\linewidth}@{}p{.15\linewidth}p{.15\linewidth}p{.15\linewidth}p{.15\linewidth}p{.15\linewidth}}
    \# & \multicolumn{2}{c}{Player} && \multicolumn{2}{c}{Game Master} \\
    \hline

    \theutterance \stepcounter{utterance}  
    & & & \multicolumn{4}{p{0.6\linewidth}}{
        \cellcolor[rgb]{0.9,0.9,0.9}{
            \makecell[{{p{\linewidth}}}]{
                \texttt{\tiny{[P1$\langle$GM]}}
                \texttt{You are an expert in causal inference. The following question is not a typical commonsense query, but rather a meticulously designed question created by a professor specializing in causal inference, intended to assess the students' mastery of the course content.} \\
\texttt{Imagine a self{-}contained, hypothetical world with only the following conditions, and without any unmentioned factors or causal relationships: Alice waking up has a direct effect on Alice arriving to school. Traffic has a direct effect on Alice arriving to school. For days when Alice wakes up on time, the probability of arriving to school on time is 46\%. For days when Alice wakes up late, the probability of arriving to school on time is 78\%. For days when Alice wakes up late, would it be more likely to see arriving to school on time if Alice had gotten up on time?} \\
\texttt{Start your answer with 'Yes' or 'No', followed by additional reasoning or evidence to support your explanation.} \\
            }
        }
    }
     \\ \\

    \theutterance \stepcounter{utterance}  
    & \multicolumn{4}{p{0.6\linewidth}}{
        \cellcolor[rgb]{0.561,0.737,0.561}{
            \makecell[{{p{\linewidth}}}]{
                \texttt{\tiny{[P1$\rangle$GM]}}
                \texttt{No} \\
\\ 
\texttt{**Reasoning:**} \\
\\ 
\texttt{To determine the answer, we must analyze the causal structure and the specific probabilities provided within the hypothetical world.} \\
\\ 
\texttt{1.  **Causal Structure**: The problem states that "Alice waking up" has a direct effect on "Alice arriving to school." This establishes a causal link where the state of waking up (On Time vs. Late) is the cause (\$X\$), and the state of arriving (On Time vs. Late) is the effect (\$Y\$).} \\
\texttt{2.  **Counterfactual Definition**: The question asks: "For days when Alice wakes up late, would it be more likely to see arriving to school on time if Alice had gotten up on time?"} \\
\texttt{    *   Let \$X\_\{late\}\$ denote the actual state (waking up late).} \\
\texttt{    *   Let \$Y\_\{obs\}\$ denote the observed outcome given \$X\_\{late\}\$.} \\
\texttt{    *   Let \$Y\_\{on\textbackslash{}\_time\}\$ denote the potential outcome if Alice had gotten up on time (\$X\_\{on\textbackslash{}\_time\}\$).} \\
\texttt{    *   The question asks if \$P(Y\_\{on\textbackslash{}\_time\} = \textbackslash{}text\{On Time\}) > P(Y\_\{obs\} = \textbackslash{}text\{On Time\})\$.} \\
\texttt{3.  **Data Analysis**:} \\
\texttt{    *   The probability of arriving on time given Alice wakes up late is explicitly stated as **78\%**. So, \$P(Y\_\{obs\} = \textbackslash{}text\{On Time\} | X\_\{late\}) = 0.78\$.} \\
\texttt{    *   The probability of arriving on time given Alice wakes up on time is explicitly stated as **46\%**. So, \$P(Y\_\{on\textbackslash{}\_time\} = \textbackslash{}text\{On Time\} | X\_\{on\textbackslash{}\_time\}) = 0.46\$.} \\
\texttt{4.  **Comparison**:} \\
\texttt{    *   We are comparing the probability of the counterfactual outcome (\$0.46\$) against the probability of the observed outcome (\$0.78\$).} \\
\texttt{    *   \$0.46 < 0.78\$.} \\
\texttt{5.  **Conclusion**: In this specific hypothetical world, the data indicates a negative correlation between waking up on time and arriving on time. When Alice wakes up late, she actually has a higher chance (78\%) of arriving on time compared to when she wakes up on time (46\%). Therefore, if Alice had gotten up on time, she would have been *less* likely to arrive on time than she was in the actual scenario where she woke up late.} \\
\\ 
\texttt{Thus, the answer is **No**.} \\
            }
        }
    }
    & & \\ \\

    \theutterance \stepcounter{utterance}  
    & & & \multicolumn{2}{p{0.3\linewidth}}{
        \cellcolor[rgb]{0.9,0.9,0.9}{
            \makecell[{{p{\linewidth}}}]{
                \texttt{\tiny{[GM$|$GM]}}
                \texttt{No} \\
            }
        }
    }
    & & \\ \\

    \theutterance \stepcounter{utterance}  
    & & & \multicolumn{2}{p{0.3\linewidth}}{
        \cellcolor[rgb]{0.9,0.9,0.9}{
            \makecell[{{p{\linewidth}}}]{
                \texttt{\tiny{[GM$|$GM]}}
                \texttt{no} \\
            }
        }
    }
    & & \\ \\

    \theutterance \stepcounter{utterance}  
    & & & \multicolumn{2}{p{0.3\linewidth}}{
        \cellcolor[rgb]{0.9,0.9,0.9}{
            \makecell[{{p{\linewidth}}}]{
                \texttt{\tiny{[GM$|$GM]}}
                \texttt{game\_result = WIN} \\
            }
        }
    }
    & & \\ \\

\end{supertabular}
}

\end{document}
